\section{Conclusion}
\label{sec:conclusion}

\paraq{What problem did this chapter start from?}
gsplit's gain over data parallelism is not uniform. Depth and hidden size change both the cost of
its shuffles and the computation it saves, so the same system is faster or slower than DGL
depending on the model, the graph, and the interconnect (Section~\ref{sec:limitations}).

\paraq{What did we add to gsplit?}
We added two hybrids. The data-parallel hybrid runs the first $k$ layers data-parallel and the
rest split, and it improves on gsplit in $37$ of $40$ configurations; the pipeline hybrid overlaps
one micro-batch's shuffle with another's compute, and our implementation does not improve on
gsplit. The hybrids turn one system into a family indexed by the split point $k$, and the best $k$
changes with the configuration (Section~\ref{sec:hybrids}).

\paraq{How do we choose among the strategies?}
We built a model. The floorline prices each strategy's overhead per layer, redundant computation
for data parallelism and the shuffle for gsplit, against two measured resources, and its ridge is
the decision boundary between them. Scanning the layers from the farthest towards the seeds gives the
split point. Its inputs are one sampling pass and three measured rates
(Section~\ref{sec:roofline}).

\paraq{Is the model correct?}
We validated it against the measured sweep. The floorline's pick is the measured winner or one
layer from it in $32$ of $40$ settings, including every PCIe setting; the misses are on NVLink,
where it picks too few data-parallel layers (Section~\ref{sec:results}).

\paraq{What does the chapter establish?}
No single parallelisation strategy is the fastest across models, graphs, and interconnects, and the
strategy can be chosen from a model of overhead instead of a sweep on the target machine.
